\subsubsection{固定资源顺序下的最早顺延修复}
保持箱组、机型、路线、资源身份和顺序不变时，充电或交接时间变化可以沿有向依赖图传播。记任务$r$的原开始为$s_r^0$，同飞机前项为$a(r)$、同电池前项为$b(r)$，扰动后任务与恢复时长为$p'_r,c'_r$，则禁止提前的最早修复满足
\begin{equation}
 s'_r=\max\{s_r^0,\ s'_{a(r)}+p'_{a(r)},\ s'_{b(r)}+p'_{b(r)}+c'_{b(r)}\}.
 \label{eq:q2_maxplus_repair}
\end{equation}
不存在某类前项时省去对应项。原资源顺序来自可行日程，依赖图无环，按拓扑顺序就可完成递推。

该构造在既定顺序和只准向后移动的条件下逐分量最早。初始任务至少要等到原开始和资源可用，取最大值达到下界；若前项已采用最早时刻，任何合法后项都不能早于其释放，所以最大值也给出后项最早开始。沿图归纳即可得到结论。随后再检查箱级截止：如果最早修复已经超期，继续向后移动不能恢复及时性，应改变资源顺序或任务结构。

这使工期变化对应最长依赖路径的增长，而不是扰动总和。两条独立链各推迟30 s可能只延长30 s；同一链上的多个变化则可能累积。单任务延迟实验固定了扰动位置，因此每个恢复结果都可以追溯到具体资源前后关系。

\subsubsection{从压力曲线形成分级调整规则}
充电时间增加30\%时，顺延工期增加415.160730 s，约7.29\%。并行作业吸收了一部分能源恢复延长，只有进入关键接续的部分才传到最后返航。充电加快情景中曲线保持不变，则是因为本次修复不允许提前任务；这与重新优化条件下加快充电可能带来的收益不同。

若仅资源接续被打破而任务仍满足能量安全，优先做时刻修复；若任务自身能耗超过安全预算，改变开始时刻没有作用，应进入重新组批或机型调整。上述两个层次分别对应时间可行性与物理可行性，使压力测试结果可以转为清楚的操作规则。

20次同起点、同预算搜索提供的是解分布，而保存正式方案提供当前最好可行日程，两者在用途上互补。Dolan与Moré提出用明确的性能度量和统一比较对象评估优化软件\cite{dolan2002}。本文相应保留起点、轮数、耗时和全部样本，将短预算稳定性与历史最好解的质量认证分开。研究记录由此既呈现有效改进，也能准确说明计算预算所支持的结论。
